\clearpage
\begin{examplebox}
\textbf{\textcolor{CExample}{例一：从一小截直线，自己“画”出周期（初中生也能懂）}}
\begin{center}\begin{tikzpicture}[font=\small,>=stealth,line join=round]
\draw[gray!40,rounded corners=3pt] (0,0) rectangle (14.2,8.3);\node at (7.1,7.700000000000001) {例1：从第一段直线出发，按对称逐段拼出图像};
\begin{scope}[shift={(2.7666666666666666,3.8842857142857143)},x=0.6666666666666666cm,y=1.3142857142857143cm]\draw[->,gray] (-2.5,0)--(6.5,0) node[right]{$x$};\draw[->,gray] (0,-1.7)--(0,1.8) node[above]{$y$};\draw[blue,thick,solid,domain=-2:-1,samples=180] plot(\x,{-1*\x+(-2)});\draw[blue,thick,solid,domain=-1:0,samples=180] plot(\x,{1*\x+(0)});\draw[blue,thick,solid,domain=0:1,samples=180] plot(\x,{1*\x+(0)});\draw[blue,thick,solid,domain=1:2,samples=180] plot(\x,{-1*\x+(2)});\draw[blue,thick,solid,domain=2:3,samples=180] plot(\x,{-1*\x+(2)});\draw[blue,thick,solid,domain=3:4,samples=180] plot(\x,{1*\x+(-4)});\draw[blue,thick,solid,domain=4:5,samples=180] plot(\x,{1*\x+(-4)});\draw[blue,thick,solid,domain=5:6,samples=180] plot(\x,{-1*\x+(6)});\fill[orange!85!black] (0.5,0.5) circle(2.2pt);\fill[orange!85!black] (2.5,-0.5) circle(2.2pt);\fill[orange!85!black] (0,0) circle(2.2pt);\fill[orange!85!black] (1,1) circle(2.2pt);\fill[orange!85!black] (2,0) circle(2.2pt);\fill[orange!85!black] (3,-1) circle(2.2pt);\fill[orange!85!black] (4,0) circle(2.2pt);\foreach \v in {-2,-1,1,2,3,4,5,6}{\node[below=4pt] at (\v,-1.35){$\v$};}\end{scope}\node[anchor=west,align=left,text width=5.0cm] at (8.4,4.0) {$a=1$，奇函数\\[9pt]轴为 $x=1$\\[9pt]$f(\frac12)=\frac12$\\[9pt]$f(\frac52)=-\frac12$\\[9pt]$f(2026.5)=-\frac12$};
\node[text width=12.6cm,align=center] at (7.1,.65) {每段长度1；0到1先上升，1到2下降，再用反号关系补齐2到4。};
\end{tikzpicture}\end{center}



\textbf{题目：}一个函数是奇函数，图像关于 $x=1$ 对称。在 $0\le x\le1$ 上，它就是 $f(x)=x$。画出 $0\le x\le4$ 的图像，并求 $f(2026.5)$。

\textbf{一步一步来：}
\begin{enumerate}[leftmargin=*,itemsep=0.3em]
 \item 先画 $(0,0)$ 到 $(1,1)$，这是一条上升的直线。
 \item 关于 $x=1$ 翻折，所以 $(1,1)$ 到 $(2,0)$ 是一条下降的直线；$1\le x\le2$ 时 $f(x)=2-x$。
 \item 由上一结论 $f(x+2)=-f(x)$，所以把 $[0,2]$ 上图像向右移 $2$ 格，再关于 $x$ 轴翻折，便得 $[2,4]$ 上的图像。点依次为 $(2,0)$、$(3,-1)$、$(4,0)$。
 \item 于是每隔 $4$ 重复一次：$2026.5=4\times506+2.5$，因此 $f(2026.5)=f(2.5)$。
 \item 因为 $2.5=0.5+2$，所以 $f(2.5)=-f(0.5)=-0.5$。
\end{enumerate}
\[
\boxed{f(2026.5)=-\tfrac12}.
\]
\textbf{图像抓手：}只需记住第一段从 $(0,0)$ 上升到 $(1,1)$，然后按对称规律翻折、翻折，整条图像便被画出来。上图已直接标出对应的分段折线。

\medskip

\end{examplebox}
\clearpage
\begin{examplebox}
\textbf{\textcolor{CExample}{例二：偶函数加点对称，也能只画折线}}
\begin{center}\begin{tikzpicture}[font=\small,>=stealth,line join=round]
\draw[gray!40,rounded corners=3pt] (0,0) rectangle (14.2,8.3);\node at (7.1,7.700000000000001) {例2：偶折线在0和4取1，在2处取负1};
\begin{scope}[shift={(2.7666666666666666,3.8842857142857143)},x=0.6666666666666666cm,y=1.3142857142857143cm]\draw[->,gray] (-2.5,0)--(6.5,0) node[right]{$x$};\draw[->,gray] (0,-1.7)--(0,1.8) node[above]{$y$};\draw[teal,thick,solid,domain=-2:-1,samples=180] plot(\x,{1*\x+(1)});\draw[teal,thick,solid,domain=-1:0,samples=180] plot(\x,{1*\x+(1)});\draw[teal,thick,solid,domain=0:1,samples=180] plot(\x,{-1*\x+(1)});\draw[teal,thick,solid,domain=1:2,samples=180] plot(\x,{-1*\x+(1)});\draw[teal,thick,solid,domain=2:3,samples=180] plot(\x,{1*\x+(-3)});\draw[teal,thick,solid,domain=3:4,samples=180] plot(\x,{1*\x+(-3)});\draw[teal,thick,solid,domain=4:5,samples=180] plot(\x,{-1*\x+(5)});\draw[teal,thick,solid,domain=5:6,samples=180] plot(\x,{-1*\x+(5)});\fill[orange!85!black] (0,1) circle(2.2pt);\fill[orange!85!black] (1,0) circle(2.2pt);\fill[orange!85!black] (2,-1) circle(2.2pt);\fill[orange!85!black] (3,0) circle(2.2pt);\fill[orange!85!black] (4,1) circle(2.2pt);\foreach \v in {-2,-1,1,2,3,4,5,6}{\node[below=4pt] at (\v,-1.35){$\v$};}\end{scope}\node[anchor=west,align=left,text width=5.0cm] at (8.4,4.0) {$g$ 为偶函数\\[9pt]中心 $(1,0)$\\[9pt]$[0,2]$ 上 $g=1-x$\\[9pt]$g(2)=-1$\\[9pt]$g(2026)=g(2)=-1$};
\node[text width=12.6cm,align=center] at (7.1,.65) {关于y轴等高，关于点(1,0)相反；两种对称一起使得隔2变号。};
\end{tikzpicture}\end{center}



\textbf{题目：}偶函数 $g(x)$ 的图像关于点 $(1,0)$ 对称，且在 $0\le x\le2$ 上有 $g(x)=1-x$。求 $g(2026)$。

\textbf{画图：}先画从 $(0,1)$ 经过 $(1,0)$ 到 $(2,-1)$ 的直线。由 $g(x+2)=-g(x)$，在 $[2,4]$ 上画从 $(2,-1)$ 经过 $(3,0)$ 到 $(4,1)$ 的直线。可以看到 $0$ 和 $4$ 高度相同，$1$ 和 $3$ 都在 $x$ 轴上。

\textbf{计算：}$2026=4\times506+2$，因此
\[
g(2026)=g(2)=1-2=\boxed{-1}.
\]
\textbf{为什么这里 $g(1)=0$？}因为 $(1,0)$ 是图像的对称中心，对称中心自身的高度只能是 $0$。这和上一种搭配的 $f(0)=0$ 是两个不同的零点来源。

\medskip

\end{examplebox}
\clearpage
\begin{examplebox}
\textbf{\textcolor{CExample}{例三：没有具体式子，也能算很大的函数值}}
\begin{center}\begin{tikzpicture}[font=\small,>=stealth,line join=round]
\draw[gray!40,rounded corners=3pt] (0,0) rectangle (14.2,8.3);\node at (7.1,7.700000000000001) {例3：保持偶函数和中心不变，把初值取为3};
\begin{scope}[shift={(1.861904761904762,3.8706896551724137)},x=0.9523809523809524cm,y=0.5287356321839081cm]\draw[->,gray] (-0.8,0)--(5.5,0) node[right]{$x$};\draw[->,gray] (0,-4.2)--(0,4.5) node[above]{$y$};\draw[teal,thick,solid,domain=-0.4:5,samples=180] plot(\x,{3*cos(deg(pi*\x/2))});\fill[orange!85!black] (0,3) circle(2.2pt);\fill[orange!85!black] (2,-3) circle(2.2pt);\fill[orange!85!black] (4,3) circle(2.2pt);\fill[orange!85!black] (1,0) circle(2.2pt);\end{scope}\node[anchor=west,align=left,text width=5.0cm] at (8.4,4.0) {$f=3\cos\frac{\pi x}2$\\[9pt]偶函数，中心 $(1,0)$\\[9pt]$f(0)=f(4)=3$\\[9pt]$2024=506\times4$\\[9pt]$f(2024)=3$};
\node[text width=12.6cm,align=center] at (7.1,.65) {图中的初值3与本题一致；不要求先猜出这条示例曲线才能求值。};
\end{tikzpicture}\end{center}



\textbf{题目：}设 $f(x)$ 是定义在 $\mathbb{R}$ 上的偶函数，图像关于 $(1,0)$ 对称，$f(0)=3$，求 $f(2024)$。

由偶函数、点对称两条信息，直接得 $f(x+2)=-f(x)$，继而 $f(x+4)=f(x)$。因为
\[
2024=4\times506+0,
\]
所以 $f(2024)=f(0)=\boxed{3}$。\textbf{题目没有让求最小正周期，完全不必去猜它。}

\medskip

\end{examplebox}
\clearpage
\begin{examplebox}
\textbf{\textcolor{CExample}{例四：高考常见的“先化小，再用奇偶性”}}
\begin{center}\begin{tikzpicture}[font=\small,>=stealth,line join=round]
\draw[gray!40,rounded corners=3pt] (0,0) rectangle (14.2,8.3);\node at (7.1,7.700000000000001) {例4：基本段是指数曲线，翻折后仍隔4重复};
\begin{scope}[shift={(1.7923076923076926,3.8842857142857143)},x=1.153846153846154cm,y=1.3142857142857143cm]\draw[->,gray] (-0.6,0)--(4.6,0) node[right]{$x$};\draw[->,gray] (0,-1.7)--(0,1.8) node[above]{$y$};\draw[blue,thick,solid,domain=0:1,samples=180] plot(\x,{2^(\x)-1});\draw[blue,thick,solid,domain=1:2,samples=180] plot(\x,{2^(2-\x)-1});\draw[blue,thick,solid,domain=2:3,samples=180] plot(\x,{1-2^(\x-2)});\draw[blue,thick,solid,domain=3:4,samples=180] plot(\x,{1-2^(4-\x)});\fill[orange!85!black] (0,0) circle(2.2pt);\fill[orange!85!black] (1,1) circle(2.2pt);\fill[orange!85!black] (2,0) circle(2.2pt);\fill[orange!85!black] (3,-1) circle(2.2pt);\fill[orange!85!black] (4,0) circle(2.2pt);\foreach \v in {1,2,3,4}{\node[below=4pt] at (\v,-1.35){$\v$};}\end{scope}\node[anchor=west,align=left,text width=5.0cm] at (8.4,4.0) {$[0,1]$ 上 $f=2^x-1$\\[9pt]奇函数，轴 $x=1$\\[9pt]$f(1)=1,\ f(3)=-1$\\[9pt]$2023=505\times4+3$\\[9pt]$f(2023)=-1$};
\node[text width=12.6cm,align=center] at (7.1,.65) {仅画0到4的部分；将它按周期4延拓到全体实数，得到符合题设的函数。};
\end{tikzpicture}\end{center}



\textbf{题目：}奇函数 $f$ 满足 $f(1+t)=f(1-t)$，且 $0\le x\le1$ 时 $f(x)=2^x-1$。求 $f(2023)$。

\textbf{第一步：}奇函数 $+$ $x=1$ 对称，故 $f(x+2)=-f(x)$，$f(x+4)=f(x)$。

\textbf{第二步：}把大数变小：
\[
2023=4\times505+3\quad\Rightarrow\quad f(2023)=f(3).
\]
\textbf{第三步：}用 $f(x+2)=-f(x)$ 更快：
\[
f(3)=-f(1)=-\bigl(2^1-1\bigr)=\boxed{-1}.
\]
\textbf{更快的诀窍：}当余数是 $3$ 时，与其往前退 $4$，不如往前退 $2$，只需变一次负号。

\medskip

\end{examplebox}
\clearpage
\begin{examplebox}
\textbf{\textcolor{CExample}{例五：专门识破“最小正周期”陷阱}}
\begin{center}\begin{tikzpicture}[font=\small,>=stealth,line join=round]
\draw[gray!40,rounded corners=3pt] (0,0) rectangle (14.2,8.8);\node at (7.1,8.200000000000001) {例5：同样奇函数、同一条轴，最小周期仍可不同};
\node at (3.1500000000000004,7.3) {$p=\sin\frac{\pi x}2$};\begin{scope}[shift={(1.19,4.338709677419355)},x=0.9800000000000001cm,y=1.225806451612903cm]\draw[->,gray] (-0.5,0)--(4.5,0) node[right]{$x$};\draw[->,gray] (0,-1.5)--(0,1.6) node[above]{$y$};\draw[blue,thick,solid,domain=-0.3:4.3,samples=180] plot(\x,{sin(deg(pi*\x/2))});\fill[orange!85!black] (1,1) circle(2pt);\fill[orange!85!black] (3,-1) circle(2pt);\draw[violet,dashed] (1,-1.3)--(1,1.3);\end{scope}\node[align=center,text width=5.8cm] at (3.1500000000000004,1.7) {关于 $x=1$ 对称\\最小正周期4};\node at (9.95,7.3) {$q=\sin\frac{3\pi x}2$};\begin{scope}[shift={(7.99,4.338709677419355)},x=0.9800000000000001cm,y=1.225806451612903cm]\draw[->,gray] (-0.5,0)--(4.5,0) node[right]{$x$};\draw[->,gray] (0,-1.5)--(0,1.6) node[above]{$y$};\draw[violet,thick,solid,domain=-0.3:4.3,samples=180] plot(\x,{sin(deg(3*pi*\x/2))});\fill[orange!85!black] (1,-1) circle(2pt);\draw[violet,dashed] (1,-1.3)--(1,1.3);\end{scope}\node[align=center,text width=5.8cm] at (9.95,1.7) {也关于 $x=1$ 对称\\最小正周期 $\frac43$};
\node[text width=12.6cm,align=center] at (7.1,.65) {两图同尺度：相同条件保证4都是周期，却不能保证最小周期相同。};
\end{tikzpicture}\end{center}



\textbf{题目：}两个函数都为奇函数，也都关于直线 $x=1$ 对称：
\[
p(x)=\sin\frac{\pi x}{2},\qquad q(x)=\sin\frac{3\pi x}{2}.
\]
它们的最小正周期相同吗？

\textbf{第一步：}两者都是奇函数，因为 $\sin(-u)=-\sin u$。

\textbf{第二步：}设 $t$ 为任意实数：
\[
\begin{aligned}
p(1+t)&=\sin(\tfrac\pi2+\tfrac{\pi t}{2})=\cos\tfrac{\pi t}{2}=p(1-t),\\
q(1+t)&=\sin(\tfrac{3\pi}{2}+\tfrac{3\pi t}{2})=-\cos\tfrac{3\pi t}{2}=q(1-t).
\end{aligned}
\]
所以二者都满足同一条轴对称条件，都以 $4$ 为一个周期。

\textbf{第三步：}但 $p(x)$ 中角度变化 $2\pi$ 需要 $x$ 增加 $4$；$q(x)$ 中角度变化 $2\pi$ 只需要 $x$ 增加 $\frac43$。所以
\[
\boxed{T_{p,\min}=4,\qquad T_{q,\min}=\tfrac43}.
\]
\textbf{这就说明：}同样的奇偶性、同样的对称轴，并不等于同样的最小正周期。这是原结论中最值得特别修正的地方。

\medskip

\end{examplebox}
\clearpage
\begin{examplebox}
\textbf{\textcolor{CExample}{例六：当 $a<0$，不要把周期写成负数}}
\begin{center}\begin{tikzpicture}[font=\small,>=stealth,line join=round]
\draw[gray!40,rounded corners=3pt] (0,0) rectangle (14.2,8.3);\node at (7.1,7.700000000000001) {例6：轴在负1，向左走4与向右走4都能重复};
\begin{scope}[shift={(4.1,3.8842857142857143)},x=0.6cm,y=1.3142857142857143cm]\draw[->,gray] (-5,0)--(5,0) node[right]{$x$};\draw[->,gray] (0,-1.7)--(0,1.8) node[above]{$y$};\draw[blue,thick,solid,domain=-4.5:4.5,samples=180] plot(\x,{sin(deg(pi*\x/2))});\fill[orange!85!black] (-1,-1) circle(2.2pt);\fill[orange!85!black] (0,0) circle(2.2pt);\fill[orange!85!black] (2,0) circle(2.2pt);\fill[orange!85!black] (4,0) circle(2.2pt);\draw[violet,dashed] (-1,-1.4)--(-1,1.4);\end{scope}\node[anchor=west,align=left,text width=5.0cm] at (8.4,4.0) {$h=\sin\frac{\pi x}2$\\[9pt]奇函数，轴 $x=-1$\\[9pt]$h(x-2)=-h(x)$\\[9pt]$h(x+4)=h(x)$\\[9pt]正周期可取4};
\node[text width=12.6cm,align=center] at (7.1,.65) {a为负时，移动方向变了，正的重复距离仍为4|a|。};
\end{tikzpicture}\end{center}



\textbf{题目：}奇函数 $h$ 关于直线 $x=-1$ 对称。能得到什么？

代 $a=-1$，先得到 $h(x-2)=-h(x)$，把 $x$ 换成 $x+2$，便有 $h(x+2)=-h(x)$。所以 $h(x+4)=h(x)$；\textbf{$4$ 是一个正周期，并不保证最小。}如果已知 $h(0)=0$，还可得 $h(2)=0$、$h(4)=0$。
\end{examplebox}
